% ============================================================
%  Chapter 4: Árboles Binarios de Búsqueda (BST)
%  Template: latex-catppuccin-academic  |  Motor: pdflatex
%
%  Dependencias heredadas del Chapter 3:
%      \usepackage{forest}   % ya debe estar en config/packages.tex
% ============================================================

\chapter{Árboles Binarios de Búsqueda (\textit{Binary Search Trees})}
\label{ch:bst}

Un árbol binario genérico, tal como se estudió en el capítulo anterior, impone únicamente
una restricción estructural: cada nodo tiene a lo sumo dos hijos. Esta restricción no dice
nada sobre el \emph{contenido} de los nodos ni sobre la relación entre las claves almacenadas
en distintas posiciones del árbol. Un \textbf{árbol binario de búsqueda} (\textit{Binary
Search Tree}, BST) agrega una \emph{propiedad de orden} que convierte la estructura en una
herramienta de búsqueda eficiente, y es precisamente la base interna de \texttt{TreeMap} y
\texttt{TreeSet} en Java.

% ------------------------------------------------------------
\section{Definición y Propiedad de Orden}
\label{sec:bst-definicion}
% ------------------------------------------------------------

\begin{definicion}{Árbol Binario de Búsqueda (BST)}
Un \textbf{BST} es un árbol binario en el que cada nodo almacena una clave única, y para
todo nodo $v$ con clave $k(v)$ se cumple:
\[
    \forall\, u \in T_{v.\mathit{left}}  : k(u) < k(v),
    \qquad
    \forall\, w \in T_{v.\mathit{right}} : k(w) > k(v),
\]
donde $T_{v.\mathit{left}}$ y $T_{v.\mathit{right}}$ denotan los subárboles izquierdo y
derecho de $v$, respectivamente. La propiedad es \textbf{recursiva}: debe satisfacerse
en cada nodo del árbol, no únicamente en la raíz.
\end{definicion}

La recursividad de la propiedad es fundamental. Un árbol puede parecer localmente correcto
en la raíz y violar el orden en un nodo interior. La figura siguiente ilustra ambos casos:

\begin{figure}[htbp]
    \centering
    \begin{minipage}{0.47\textwidth}
        \centering
        \begin{forest}
            for tree={circle, draw=CtpGreen, fill=CtpGreen!20, text=CtpText,
                       minimum size=0.9cm, inner sep=1pt, l sep=12mm, s sep=8mm,
                       font=\sffamily\small}
            [30 [15 [10][20]] [50 [40][70]]]
        \end{forest}
        \par\smallskip
        {\small\color{CtpGreen}\textbf{BST válido} — la propiedad de orden se
        cumple en todos los nodos.}
    \end{minipage}
    \hfill
    \begin{minipage}{0.47\textwidth}
        \centering
        \begin{forest}
            for tree={circle, draw=CtpRed, fill=CtpRed!15, text=CtpText,
                       minimum size=0.9cm, inner sep=1pt, l sep=12mm, s sep=8mm,
                       font=\sffamily\small}
            [30 [15 [10][20]] [50 [55][70]]]
        \end{forest}
        \par\smallskip
        {\small\color{CtpRed}\textbf{No válido} — el nodo $55$ está en el subárbol
        izquierdo de $50$, pero $55 > 50$.}
    \end{minipage}
    \caption{BST válido (izquierda) vs.\ no válido (derecha). La violación en el árbol
             de la derecha ocurre en un nodo interior ($50$), no en la raíz, lo que
             ilustra la necesidad de verificar la propiedad de forma recursiva.}
    \label{fig:bst-valido-invalido}
\end{figure}

\begin{observacion}{In-order como verificación}
El recorrido in-order de un BST produce las claves en orden estrictamente ascendente. Esta
propiedad es consecuencia directa de la definición: en cada nodo $v$, el subárbol izquierdo
contiene claves menores que $k(v)$ y el subárbol derecho las mayores, que es exactamente
el orden en que los visita el in-order $(L, \text{Raíz}, R)$. Verificar que el in-order
produce una secuencia ascendente es condición necesaria y suficiente para validar que un
árbol es un BST correcto.
\end{observacion}

% ------------------------------------------------------------
\section{Búsqueda}
\label{sec:bst-busqueda}
% ------------------------------------------------------------

La búsqueda en un BST explota la propiedad de orden para descartar un subárbol completo
en cada comparación, de manera análoga a la búsqueda binaria en un arreglo ordenado.

\begin{algorithm}[H]
\caption{Búsqueda en un BST}
\label{alg:bst-search}
\KwIn{Nodo raíz $v$, clave $k$ buscada}
\KwOut{Nodo con clave $k$, o \texttt{null} si no existe}
\If{$v = \texttt{null}$}{\Return \texttt{null}}
\If{$k = k(v)$}{\Return $v$}
\If{$k < k(v)$}{\Return \texttt{search}$(v.\mathit{left},\; k)$}
\Return \texttt{search}$(v.\mathit{right},\; k)$\;
\end{algorithm}

Cada llamada recursiva descarta el subárbol que no puede contener la clave. La complejidad
es $O(H)$, donde $H$ es la altura del árbol. En un BST balanceado $H = O(\log n)$; en el
peor caso $H = O(n)$, como se verá en la \cref{sec:bst-complejidad}.

\begin{observacion}{Analogía con búsqueda binaria}
La búsqueda en un BST y la búsqueda binaria en un arreglo ordenado son estructuralmente
idénticas: en cada paso se compara la clave buscada con un elemento pivote y se descarta
la mitad del espacio de búsqueda restante. La diferencia es que en el arreglo el pivote
es el elemento central (acceso por índice), mientras que en el BST el pivote es la clave
del nodo actual y el ``descarte'' ocurre siguiendo un puntero. Ambos algoritmos tienen
complejidad $O(\log n)$ cuando la estructura está balanceada.
\end{observacion}

% ------------------------------------------------------------
\section{Inserción}
\label{sec:bst-insercion}
% ------------------------------------------------------------

La inserción sigue exactamente el mismo camino que la búsqueda. Se recorre el árbol
comparando la nueva clave con la clave de cada nodo hasta encontrar una posición
\texttt{null}, que es donde se inserta el nuevo nodo. Si la clave ya existe, no se
realiza ninguna acción (el BST no admite duplicados).

\begin{algorithm}[H]
\caption{Inserción en un BST}
\label{alg:bst-insert}
\KwIn{Nodo raíz $v$, clave $k$ a insertar}
\KwOut{Raíz del árbol actualizado}
\If{$v = \texttt{null}$}{\Return \textbf{nuevoNodo}$(k)$}
\If{$k < k(v)$}{$v.\mathit{left} \leftarrow$ \texttt{insert}$(v.\mathit{left},\; k)$}
\ElseIf{$k > k(v)$}{$v.\mathit{right} \leftarrow$ \texttt{insert}$(v.\mathit{right},\; k)$}
\Return $v$\;
\end{algorithm}

\begin{ejemplo}{Construcción de un BST por inserción secuencial}
Insertando las claves $30,\; 15,\; 50,\; 10,\; 20,\; 40,\; 70$ en ese orden en un BST
vacío, cada clave sigue el camino que usaría la búsqueda hasta encontrar la primera
posición \texttt{null}. La \cref{tab:bst-insercion} detalla la trayectoria de cada clave
y la \cref{fig:bst-insercion} muestra el árbol resultante, cuyo in-order produce
$10, 15, 20, 30, 40, 50, 70$ —secuencia estrictamente ascendente, confirmando que la
propiedad BST se preservó en todas las inserciones.
\end{ejemplo}

\begin{acadtable}[htbp]{tab:bst-insercion}{Trayectoria de inserción para cada clave}
    \begin{tabularx}{\textwidth}{c l X}
        \headerrow
        \textbf{Clave} & \textbf{Camino recorrido} & \textbf{Posición final} \\
        \hline
        $30$ & (árbol vacío)                         & Raíz \\
        $15$ & $30 \to$ izq                          & Hijo izquierdo de $30$ \\
        $50$ & $30 \to$ der                          & Hijo derecho de $30$ \\
        $10$ & $30 \to$ izq $\to 15 \to$ izq        & Hijo izquierdo de $15$ \\
        $20$ & $30 \to$ izq $\to 15 \to$ der        & Hijo derecho de $15$ \\
        $40$ & $30 \to$ der $\to 50 \to$ izq        & Hijo izquierdo de $50$ \\
        $70$ & $30 \to$ der $\to 50 \to$ der        & Hijo derecho de $50$ \\
    \end{tabularx}
\end{acadtable}

\begin{figure}[htbp]
    \centering
    \begin{forest}
        for tree={
            circle, draw=CtpOverlay2, fill=CtpSurface0, text=CtpText,
            minimum size=0.9cm, inner sep=2pt, l sep=13mm, s sep=9mm,
            font=\sffamily\small
        }
        [30 [15 [10][20]] [50 [40][70]]]
    \end{forest}
    \caption{BST resultante de insertar $30, 15, 50, 10, 20, 40, 70$ en ese orden.
             La raíz $30$ es la mediana del conjunto, lo que produce un árbol
             perfectamente balanceado de altura $2$.}
    \label{fig:bst-insercion}
\end{figure}

\begin{observacion}{El orden de inserción determina la forma del árbol}
La misma secuencia de claves produce distintos árboles según el orden en que se insertan.
El árbol del ejemplo anterior es perfectamente balanceado porque $30$ es la mediana del
conjunto. Si la primera clave insertada fuera $10$ (el mínimo), el árbol degeneraría hacia
la derecha. Este fenómeno se estudia en detalle en la \cref{sec:bst-complejidad}.
\end{observacion}

% ------------------------------------------------------------
\section{Eliminación}
\label{sec:bst-eliminacion}
% ------------------------------------------------------------

La eliminación es la operación más compleja del BST. El procedimiento depende del número
de hijos del nodo a eliminar y se divide en tres casos mutuamente excluyentes.

\subsection{Caso 1 — Nodo hoja}
\label{subsec:caso1}

\begin{definicion}{Eliminación: Caso 1}
Si el nodo a eliminar es una \textbf{hoja} (sin hijos), se elimina directamente. El padre
actualiza el puntero correspondiente a \texttt{null}.
\end{definicion}

\begin{figure}[htbp]
    \centering
    \begin{minipage}{0.45\textwidth}
        \centering
        \begin{forest}
            for tree={circle, draw=CtpOverlay2, fill=CtpSurface0, text=CtpText,
                       minimum size=0.85cm, inner sep=1pt, l sep=12mm, s sep=8mm,
                       font=\sffamily\small}
            [30 [15 [10, fill=CtpRed!30, draw=CtpRed][20]] [50]]
        \end{forest}
        \par\smallskip{\small Antes: \texttt{remove(10)}}
    \end{minipage}
    \hfill
    \begin{minipage}{0.45\textwidth}
        \centering
        \begin{forest}
            for tree={circle, draw=CtpOverlay2, fill=CtpSurface0, text=CtpText,
                       minimum size=0.85cm, inner sep=1pt, l sep=12mm, s sep=8mm,
                       font=\sffamily\small}
            [30 [15 [,phantom,draw=none,fill=none][20]] [50]]
        \end{forest}
        \par\smallskip{\small Después: el padre 15 apunta a \texttt{null}.}
    \end{minipage}
    \caption{Eliminación Caso 1: el nodo $10$ es hoja y se elimina directamente.}
    \label{fig:bst-caso1}
\end{figure}

\subsection{Caso 2 — Nodo con un hijo}
\label{subsec:caso2}

\begin{definicion}{Eliminación: Caso 2}
Si el nodo a eliminar tiene \textbf{exactamente un hijo}, el padre del nodo eliminado
adopta directamente a ese hijo: el puntero del padre que apuntaba al nodo eliminado
pasa a apuntar al único hijo del nodo eliminado.
\end{definicion}

\begin{figure}[htbp]
    \centering
    \begin{minipage}{0.45\textwidth}
        \centering
        \begin{forest}
            for tree={circle, draw=CtpOverlay2, fill=CtpSurface0, text=CtpText,
                       minimum size=0.85cm, inner sep=1pt, l sep=12mm, s sep=8mm,
                       font=\sffamily\small}
            [30 [15, fill=CtpRed!30, draw=CtpRed [,phantom,draw=none,fill=none][20]]
                [50]]
        \end{forest}
        \par\smallskip{\small Antes: \texttt{remove(15)}, que tiene solo hijo derecho.}
    \end{minipage}
    \hfill
    \begin{minipage}{0.45\textwidth}
        \centering
        \begin{forest}
            for tree={circle, draw=CtpOverlay2, fill=CtpSurface0, text=CtpText,
                       minimum size=0.85cm, inner sep=1pt, l sep=12mm, s sep=8mm,
                       font=\sffamily\small}
            [30 [20][50]]
        \end{forest}
        \par\smallskip{\small Después: el padre $30$ adopta directamente a $20$.}
    \end{minipage}
    \caption{Eliminación Caso 2: el nodo $15$ tiene un único hijo ($20$). El padre $30$
             pasa a apuntar directamente a $20$.}
    \label{fig:bst-caso2}
\end{figure}

\subsection{Caso 3 — Nodo con dos hijos}
\label{subsec:caso3}

\begin{definicion}{Eliminación: Caso 3}
Si el nodo a eliminar tiene \textbf{dos hijos}, se reemplaza su clave por la del
\textbf{predecesor in-order}: el mayor de todos sus descendientes izquierdos. Luego
se elimina el predecesor in-order de su posición original.
\end{definicion}

El predecesor in-order se localiza en $O(H)$ sin recorrer el árbol completo:

\begin{center}
\textit{Desde el nodo a eliminar: un paso a la izquierda, luego todo lo que se pueda
a la derecha.}
\end{center}

La elección del predecesor in-order como sustituto no es arbitraria. Es el único nodo
que satisface simultáneamente las dos condiciones que preservan la propiedad BST:

\begin{itemize}
    \item Es \textbf{menor que todo el subárbol derecho} del nodo eliminado, porque
          pertenece al subárbol izquierdo.
    \item Es \textbf{mayor que todo el resto del subárbol izquierdo}, porque es el
          máximo de ese subárbol.
\end{itemize}

Además, el predecesor in-order \textbf{nunca tiene hijo derecho} —si lo tuviera, ese
hijo sería mayor y él no sería el máximo del subárbol izquierdo—, por lo que su
eliminación de la posición original es siempre Caso 1 o Caso 2, nunca Caso 3.

\begin{figure}[htbp]
    \centering
    \begin{minipage}{0.45\textwidth}
        \centering
        \begin{forest}
            for tree={circle, draw=CtpOverlay2, fill=CtpSurface0, text=CtpText,
                       minimum size=0.85cm, inner sep=1pt, l sep=12mm, s sep=7mm,
                       font=\sffamily\small}
            [30, fill=CtpRed!30, draw=CtpRed
                [15 [10][20, fill=CtpBlue!35, draw=CtpBlue]]
                [50 [40][70]]]
        \end{forest}
        \par\smallskip{\small Antes: \texttt{remove(30)}. Predecesor in-order
        (en azul): $20$.}
    \end{minipage}
    \hfill
    \begin{minipage}{0.45\textwidth}
        \centering
        \begin{forest}
            for tree={circle, draw=CtpOverlay2, fill=CtpSurface0, text=CtpText,
                       minimum size=0.85cm, inner sep=1pt, l sep=12mm, s sep=7mm,
                       font=\sffamily\small}
            [20 [15 [10][,phantom,draw=none,fill=none]] [50 [40][70]]]
        \end{forest}
        \par\smallskip{\small Después: $20$ reemplaza a $30$. Su posición original
        (hijo derecho de $15$) queda vacía — Caso 1.}
    \end{minipage}
    \caption{Eliminación Caso 3: \texttt{remove(30)}. El predecesor in-order ($20$) se
             obtiene yendo un paso a la izquierda ($15$) y luego todo a la derecha ($20$,
             sin hijo derecho). Reemplaza a $30$ y su eliminación original es Caso 1.}
    \label{fig:bst-caso3}
\end{figure}

\begin{algorithm}[H]
\caption{Eliminación en un BST}
\label{alg:bst-delete}
\KwIn{Nodo raíz $v$, clave $k$ a eliminar}
\KwOut{Raíz del árbol actualizado}
\If{$v = \texttt{null}$}{\Return \texttt{null}}
\If{$k < k(v)$}{
    $v.\mathit{left} \leftarrow$ \texttt{delete}$(v.\mathit{left},\; k)$\;
    \Return $v$\;
}
\If{$k > k(v)$}{
    $v.\mathit{right} \leftarrow$ \texttt{delete}$(v.\mathit{right},\; k)$\;
    \Return $v$\;
}
\tcc{Nodo encontrado — tres casos:}
\If{$v.\mathit{left} = \texttt{null}$ \textbf{and} $v.\mathit{right} = \texttt{null}$}{
    \Return \texttt{null}\; \tcc*{Caso 1: hoja}
}
\If{$v.\mathit{left} = \texttt{null}$}{
    \Return $v.\mathit{right}$\; \tcc*{Caso 2: solo hijo derecho}
}
\If{$v.\mathit{right} = \texttt{null}$}{
    \Return $v.\mathit{left}$\; \tcc*{Caso 2: solo hijo izquierdo}
}
\tcc{Caso 3: dos hijos — buscar predecesor in-order}
$p \leftarrow v.\mathit{left}$\;
\While{$p.\mathit{right} \neq \texttt{null}$}{
    $p \leftarrow p.\mathit{right}$\;
}
$k(v) \leftarrow k(p)$\;
$v.\mathit{left} \leftarrow$ \texttt{delete}$(v.\mathit{left},\; k(p))$\;
\Return $v$\;
\end{algorithm}

% ------------------------------------------------------------
\section{Complejidad}
\label{sec:bst-complejidad}
% ------------------------------------------------------------

Toda operación del BST —búsqueda, inserción, eliminación— es $O(H)$, donde $H$ es la
altura del árbol. La complejidad efectiva depende de cuánto vale $H$, que a su vez depende
del orden en que se insertan las claves.

\begin{acadtable}[htbp]{tab:bst-complejidad}{Complejidad de las operaciones en un BST}
    \begin{tabularx}{\textwidth}{l c c}
        \headerrow
        \textbf{Operación} & \textbf{Caso promedio} & \textbf{Peor caso} \\
        \hline
        Búsqueda   & $O(\log n)$ & $O(n)$ \\
        Inserción  & $O(\log n)$ & $O(n)$ \\
        Eliminación & $O(\log n)$ & $O(n)$ \\
    \end{tabularx}
\end{acadtable}

El peor caso ocurre cuando las claves se insertan en orden estrictamente creciente o
decreciente. En ese escenario, cada nueva clave es siempre mayor (o siempre menor) que
todas las anteriores, por lo que se inserta sistemáticamente como hijo derecho (o izquierdo)
del último nodo. El árbol degenera en una lista enlazada con $H = n - 1$:

\begin{figure}[htbp]
    \centering
    \begin{minipage}{0.35\textwidth}
        \centering
        \begin{forest}
            for tree={circle, draw=CtpRed, fill=CtpRed!15, text=CtpText,
                       minimum size=0.85cm, inner sep=1pt,
                       l sep=11mm, s sep=8mm,
                       font=\sffamily\small,
                       grow=south}
            [10 [,phantom,draw=none,fill=none]
                [20 [,phantom,draw=none,fill=none]
                    [30 [,phantom,draw=none,fill=none]
                        [40 [,phantom,draw=none,fill=none]
                            [50]]]]]
        \end{forest}
    \end{minipage}
    \hfill
    \begin{minipage}{0.58\textwidth}
        \small
        Secuencia de inserción: $10, 20, 30, 40, 50$.\\[6pt]
        Cada clave es mayor que todas las anteriores,
        por lo que se inserta siempre como hijo derecho
        del nodo hoja actual.\\[6pt]
        Resultado: $H = 4 = n - 1$.\\[6pt]
        Una búsqueda de $50$ requiere $5$ comparaciones
        —equivalente a una búsqueda lineal en una lista
        de $5$ elementos.
    \end{minipage}
    \caption{Árbol BST degenerado por inserción en orden creciente. La altura es $n-1$
             en lugar del óptimo $\lfloor\log_2 n\rfloor = 2$. El árbol ha perdido toda
             ventaja sobre una lista enlazada.}
    \label{fig:bst-degenerado}
\end{figure}

\begin{observacion}{La mediana como primer elemento}
Si se elige como primera clave a insertar la mediana del conjunto, y se repite la estrategia
recursivamente en cada subárbol, el BST resultante es perfectamente balanceado con
$H = \lfloor\log_2 n\rfloor$. En la práctica esto raramente se conoce de antemano, lo que
motiva el uso de estructuras auto-balanceadas como los árboles AVL o rojo-negro, que
garantizan $H = O(\log n)$ independientemente del orden de inserción.
\end{observacion}

% ------------------------------------------------------------
\section{Ejercicios}
\label{sec:bst-ejercicios}
% ------------------------------------------------------------

\begin{ejercicio}
Inserte las siguientes claves en un BST vacío, en el orden indicado:
\[
    45,\; 25,\; 60,\; 10,\; 30,\; 55,\; 70,\; 5,\; 28,\; 50
\]
\textbf{(1)} Dibuje el árbol resultante. \textbf{(2)} Determine si es completo, lleno
y balanceado. \textbf{(3)} Escriba los recorridos pre-order, in-order y post-order.
\textbf{(4)} Verifique que el in-order produce una secuencia ascendente.
\end{ejercicio}

\begin{ejercicio}
Dado el árbol del ejercicio anterior, realice las siguientes eliminaciones en secuencia
(cada una parte del resultado de la anterior). Para cada eliminación, indique el caso
que aplica (hoja / un hijo / dos hijos), muestre el camino al predecesor in-order si
corresponde, y dibuje el árbol resultante.

\[
    \texttt{remove}(10), \quad \texttt{remove}(60), \quad \texttt{remove}(45)
\]
\end{ejercicio}

\begin{ejercicio}
Considere las dos secuencias de inserción siguientes sobre el mismo conjunto de claves
$\{5, 10, 15, 20, 25, 30\}$:
\begin{itemize}
    \item Secuencia A: $15,\; 10,\; 25,\; 5,\; 20,\; 30$
    \item Secuencia B: $5,\; 10,\; 15,\; 20,\; 25,\; 30$
\end{itemize}
\textbf{(1)} Dibuje el BST resultante para cada secuencia. \textbf{(2)} Calcule la
altura de cada árbol. \textbf{(3)} Determine cuántas comparaciones requiere buscar la
clave $30$ en cada árbol. \textbf{(4)} Explique por qué ambos árboles producen el mismo
in-order pero tienen rendimientos de búsqueda distintos.
\end{ejercicio}


\printbibliography[heading=subbibliography]
